The expected contribution of the complete work is a software architecture that improves the analysis and aggregation stage of scenario-based simulation outputs by combining distributed storage, distributed computing, and analyst-defined metric semantics. If successful, the resulting system should reduce the need for manual result transfer, shorten the time required to obtain analytical products from simulation batches, and provide a more scalable basis for future work in decision-support workflows built on \gls{asa}, helping \gls{cbp} analysis as a whole for the Brazilian Air Force.

\section{Research and Validation}

\begin{itemize}
    \item Evaluate MapReduce viability as a processing paradigm for this work's needs;
    \item Assess \gls{hdfs} for storage capability, data locality, and batch-level post-processing support;
    \item Study the existing \gls{asa} codebase and workflow to identify integration points.
\end{itemize}

\section{Implementation}

\begin{itemize}
    \item Create a user \gls{api} for defining or selecting aggregate functions before processing;
    \item Implement a \gls{dfs} for simulation outputs;
    \item Build distributed computing logic to extract metrics and compute aggregate values;
    \item Test the architecture on a local computer cluster;
    \item Progressively integrate this architecture into \gls{asa} systems.
\end{itemize}

\section{Evaluation}

\begin{itemize}
    \item Compare end-to-end processing time against the current serial pipeline for representative batches;
    \item Quantify performance gains from parallel metric extraction, reduced data transfer, and storage access improvements;
    \item Assess qualitative aspects: ease of use, deployment complexity, and coupling with the existing infrastructure;
\end{itemize}
